problem:  
Let \( (X,d) \) be a compact \(n\)-dimensional Alexandrov space of curvature \(\ge 0\). Suppose that for every \(x \in X\), the tangent cone \(T_x X\) exists and is unique, and moreover the assignment  
\[
x \mapsto T_x X
\]
is Hölder continuous in the pointed Gromov–Hausdorff topology.  

**Question:**  
Does such Hölder continuity of tangent cones imply that the regular part \(M_{\mathrm{reg}}\subset X\) admits a \(C^{1,\alpha}\) Riemannian metric structure compatible with the Alexandrov metric, and that \(X\) itself is homeomorphic to a smooth manifold (perhaps after neglecting a set of Hausdorff codimension \(\ge 3\))?  

In other words, can quantitative continuity of tangent cones upgrade the canonical Lipschitz manifold structure of \(X\) to a \(C^{1,\alpha}\)-smooth structure compatible with a nonnegatively curved Riemannian geometry?

Why is it a "good" problem:  
1. **Conceptual depth:**  
   Alexandrov spaces already possess tangent cones and a stratification by the types of their spaces of directions, but known regularity ends essentially at *topological* or *Lipschitz* smoothness. This problem asks whether an additional analytic assumption—quantitative continuity of tangent cones—forces a genuine differentiable or Riemannian structure. It thus lies at the forefront of efforts to understand when metric curvature bounds yield smooth differentiable geometry.

2. **Bridge between fields:**  
   The question connects **metric geometry** (Gromov–Hausdorff convergence, Alexandrov curvature), **geometric analysis** (quantitative regularity and Hölder continuity), and **differential topology** (existence of smooth structures). It sits precisely at the crossroad where comparison geometry might transition into smooth Riemannian geometry, and would illuminate the boundary between purely metric and smooth curvature worlds.

3. **Potential significance:**  
   A positive resolution would establish a powerful new *regularity principle*: sufficiently controlled cone behavior forces differentiability and geometric tameness under nonnegative curvature. This would parallel the role that quantitative tangent-cone stability plays in the Cheeger–Colding theory for Ricci limit spaces. Conversely, a counterexample would reveal subtle new obstructions to smoothing Alexandrov spaces, marking fundamental limitations of metric regularity.

4. **Simplicity and profundity:**  
   The statement is minimal and purely geometric: “Does Hölder continuity of tangent cones smooth an Alexandrov space?” Its resolution would require deep synthesis of metric analysis, curvature bounds, and topological regularity theory. The question’s elementary formulation but formidable depth exemplifies the “narrow entrance, profound depth” ideal of a truly good mathematical problem.